\input cwebmac

\N{1}{1}Swim Times.

This program fetches best short-course yard (SCY) times for two College
Area Swim Team (CAST) swimmers:
\medskip
$${\vbox{
{\bf Stella Julianna Evans} --- 10~\&~Under Girls

{\bf Kalea Benavente} --- 13--14 Girls
}}$$

\medskip\noindent
Twelve events are reported for each swimmer:
50, 100, 200, and 500~Freestyle; 50 and 100~Butterfly;
50 and 100~Backstroke; 50 and 100~Breaststroke; and 100 and 200~Individual
Medley.  Each time record includes the swim date and motivational standard
attained (B, BB, A, \dots).  Data is fetched live from the USA~Swimming
data hub, which is powered by the Sisense analytics platform.

\fi

\M{2}{\bf How it works.}  The USA~Swimming data hub exposes a Sisense
JAQL~API.  We make two kinds of requests:

\medskip\item{1.} A {\it person search\/} to resolve the swimmer's internal
\PB{\\{PersonKey}}, given a search string and a substring to match the full
name.

\item{2.} A {\it times query\/} for each event, filtered by \PB{\\{PersonKey}}
and event code (e.g.\ \.{100 FR SCY}).  Each row delivers five cells:
swim time, sort key, meet name, swim date (as \.{YYYYMMDD}), and
motivational standard.
\medskip

All HTTP communication is handled by \.{libcurl}.  JSON responses are
scanned with simple string operations rather than a full parse tree.

\fi

\M{3}{\bf Compilation.}  After tangling with \.{ctangle}:
$$\.{gcc -O2 -o swim-times swim-times.c \$(curl-config --libs)}$$

\fi

\M{4}{\bf Program structure.}  The top-level arrangement of the tangled
output is:

\Y\B\X6:Includes\X\X7:Global constants\X\X9:Type definitions\X\X12:HTTP
utilities\X\X16:JSON scanner\X\X19:Person lookup\X\X21:Times fetch\X\X25:Main
function\X\par
\fi

\N{1}{5}Includes and constants.

\fi

\M{6}\B\X6:Includes\X${}\E{}$\6
\8\#\&{include} \.{<stdio.h>}\6
\8\#\&{include} \.{<stdlib.h>}\6
\8\#\&{include} \.{<string.h>}\6
\8\#\&{include} \.{<curl/curl.h>}\par
\U4.\fi

\M{7}The Sisense bearer token authenticates us against the analytics
instance that backs \.{data.usaswimming.org}.  \PB{\.{EVENTS}} lists the
twelve SCY event codes we query.

\Y\B\4\X7:Global constants\X${}\E{}$\6
\8\#\&{define} \.{NUM\_EVENTS}\5\T{12}\6
\8\#\&{define} \.{MAX\_TIMES}\5\T{200}\6
\&{static} \&{const} \&{char} \.{SISENSE\_TOKEN}[\,]${}\K%
\.{"eyJhbGciOiJIUzI1NiI}\)\.{sInR5cCI6IkpXVCJ9"}\.{".eyJ1c2VyIjoiNjY0Ym}\)%
\.{E2NmE5M2ZiYTUwMDM4NW}\)\.{IyMWQwIiwiYXBpU2Vjcm}\)\.{V0Ijo"}%
\.{"iNDQ0YTE3NWQtM2I1OC}\)\.{03NDhhLTVlMGEtYTVhZD}\)\.{E2MmRmODJlIiwiYWxsb3}\)%
\.{dlZFRl"}\.{"bmFudHMiOlsiNjRhYzE}\)\.{5ZTEwZTkxNzgwMDFiYzM}\)%
\.{5YmVhIl0sInRlbmFudEl}\)\.{kIjoiN"}\.{"jRhYzE5ZTEwZTkxNzgw}\)%
\.{MDFiYzM5YmVhIn0"}\.{".izSIvaD2udKTs3QRng}\)\.{la1Aw23kZVyoq7Xh23Ab}\)%
\.{PUw1M"};{}$\6
\&{static} \&{const} \&{char} \.{SISENSE\_API}[\,]${}\K\.{"https://usaswimming}%
\)\.{.sisense.com/api/dat}\)\.{asources"};{}$\6
\&{static} \&{const} \&{char} ${}{*}\.{EVENTS}[\.{NUM\_EVENTS}]\K\{\.{"50\ FR\
SCY"},\39\.{"100\ FR\ SCY"},\39\.{"200\ FR\ SCY"},\39\.{"500\ FR\ SCY"},\39%
\.{"50\ FL\ SCY"},\39\.{"100\ FL\ SCY"},\39\.{"50\ BK\ SCY"},\39\.{"100\ BK\
SCY"},\39\.{"50\ BR\ SCY"},\39\.{"100\ BR\ SCY"},\39\.{"100\ IM\ SCY"},\39%
\.{"200\ IM\ SCY"}\}{}$;\par
\U4.\fi

\N{1}{8}Data structures.

\fi

\M{9}A \PB{\\{Buffer}} holds a dynamically-grown heap string.  \.{libcurl}
appends each response chunk to it via the write callback.

\Y\B\4\X9:Type definitions\X${}\E{}$\6
\&{typedef} \&{struct} ${}\{{}$\1\6
\&{char} ${}{*}\\{data};{}$\6
\&{size\_t} \\{size};\2\6
${}\}{}$ \&{Buffer};\par
\A10.
\U4.\fi

\M{10}A \PB{\\{TimeRow}} records one swim result.  \PB{\\{sort\_key}} is a
numeric proxy
for the swim time (smaller is faster); \PB{\\{time}} is the formatted string
(e.g.\ \.{1:02.45}); \PB{\\{date}} is the swim date formatted as
\.{YYYY-MM-DD}; \PB{\\{standard}} is the motivational standard attained
(e.g.\ \.{B}, \.{BB}, \.{A}, \.{Slower Than B}); and \PB{\\{meet}} is the
meet name.  The Sisense API returns the sort key as a quoted decimal
string such as \.{"1010007029.00"}, so we store it as a \PB{\&{double}}.

\Y\B\4\X9:Type definitions\X${}\mathrel+\E{}$\6
\&{typedef} \&{struct} ${}\{{}$\1\6
\&{double} \\{sort\_key};\6
\&{char} \\{time}[\T{32}];\6
\&{char} \\{date}[\T{16}];\6
\&{char} \\{standard}[\T{48}];\6
\&{char} \\{meet}[\T{256}];\2\6
${}\}{}$ \&{TimeRow};\par
\fi

\N{1}{11}HTTP utilities.

\fi

\M{12}The \.{libcurl} write callback appends each incoming chunk to a
\PB{\&{Buffer}}, growing the allocation as needed.  It maintains a null
terminator so the buffer can always be treated as a C string.

\Y\B\4\X12:HTTP utilities\X${}\E{}$\6
\1\1\&{static} \&{size\_t} \\{write\_cb}(\&{void} ${}{*}\\{ptr},\39{}$\&{size%
\_t} \\{size}${},\39{}$\&{size\_t} \\{nmemb}${},\39{}$\&{void} ${}{*}\\{ud})\2%
\2{}$\6
${}\{{}$\1\6
\&{Buffer} ${}{*}\|b\K{}$(\&{Buffer} ${}{*}){}$ \\{ud};\6
\&{size\_t} \|n${}\K\\{size}*\\{nmemb};{}$\6
\&{char} ${}{*}\|p\K\\{realloc}(\|b\MG\\{data},\39\|b\MG\\{size}+\|n+\T{1});{}$%
\7
\&{if} ${}(\R\|p){}$\1\5
\&{return} \T{0};\2\6
${}\|b\MG\\{data}\K\|p;{}$\6
${}\\{memcpy}(\|b\MG\\{data}+\|b\MG\\{size},\39\\{ptr},\39\|n);{}$\6
${}\|b\MG\\{size}\MRL{+{\K}}\|n;{}$\6
${}\|b\MG\\{data}[\|b\MG\\{size}]\K\.{'\\0'};{}$\6
\&{return} \|n;\6
\4${}\}{}$\2\par
\A13.
\U4.\fi

\M{13}\PB{\\{post\_json}} performs a single HTTP~POST with a JSON body and
returns
the response as a null-terminated heap string, or \PB{$\NULL$} on failure.
The caller is responsible for freeing the returned string.

\Y\B\4\X12:HTTP utilities\X${}\mathrel+\E{}$\6
\1\1\&{static} \&{char} ${}{*}{}$\\{post\_json}(\&{const} \&{char} ${}{*}%
\\{url},\39{}$\&{const} \&{char} ${}{*}\\{body})\2\2{}$\6
${}\{{}$\1\6
${}\.{CURL}*\\{curl}\K\\{curl\_easy\_init}(\,);{}$\6
\&{if} ${}(\R\\{curl}){}$\1\5
\&{return} ${}\NULL;{}$\2\7
\&{Buffer} \\{buf}${}\K\{\NULL,\39\T{0}\};{}$\6
\&{char} \\{auth\_hdr}[\T{1600}];\7
${}\\{snprintf}(\\{auth\_hdr},\39{}$\&{sizeof} \\{auth\_hdr}${},\39%
\.{"Authorization:\ Bear}\)\.{er\ \%s"},\39\.{SISENSE\_TOKEN});{}$\7
\&{struct} \\{curl\_slist} ${}{*}\\{hdrs}\K\NULL;{}$\7
${}\\{hdrs}\K\\{curl\_slist\_append}(\\{hdrs},\39\.{"Content-Type:\ appli}\)%
\.{cation/json"});{}$\6
${}\\{hdrs}\K\\{curl\_slist\_append}(\\{hdrs},\39\\{auth\_hdr});{}$\6
${}\\{curl\_easy\_setopt}(\\{curl},\39\.{CURLOPT\_URL},\39\\{url});{}$\6
${}\\{curl\_easy\_setopt}(\\{curl},\39\.{CURLOPT\_HTTPHEADER},\39\\{hdrs});{}$\6
${}\\{curl\_easy\_setopt}(\\{curl},\39\.{CURLOPT\_POSTFIELDS},\39\\{body});{}$\6
${}\\{curl\_easy\_setopt}(\\{curl},\39\.{CURLOPT\_WRITEFUNCTION},\39\\{write%
\_cb});{}$\6
${}\\{curl\_easy\_setopt}(\\{curl},\39\.{CURLOPT\_WRITEDATA},\39{\AND}%
\\{buf});{}$\6
${}\\{CURLcode}\\{rc}\K\\{curl\_easy\_perform}(\\{curl});{}$\6
\\{curl\_slist\_free\_all}(\\{hdrs});\6
\\{curl\_easy\_cleanup}(\\{curl});\6
\&{if} ${}(\\{rc}\I\.{CURLE\_OK}){}$\5
${}\{{}$\1\6
${}\\{free}(\\{buf}.\\{data});{}$\6
\&{return} ${}\NULL;{}$\6
\4${}\}{}$\2\6
\&{return} \\{buf}${}.\\{data};{}$\6
\4${}\}{}$\2\par
\fi

\N{1}{14}JSON scanner.

\fi

\M{15}Sisense JAQL responses contain a top-level \PB{\.{"values"}} array.  Each
element is itself an array of objects, each holding either a \PB{\.{"text"}}
key (string value) or a \PB{\.{"data"}} key (numeric value).  A typical
person-search row looks like:
$$\hbox{\.{[{"text":"Stella Julianna
Evans"},{"data":12345},{"text":"CAST"}]}}$$

Rather than build a full parse tree we scan forward for a key and
extract the immediately following value, advancing a position pointer.

\fi

\M{16}\PB{\\{scan\_string}} finds the first \PB{\.{"key"}} at or after \PB{${*}%
\\{pos}$}, copies the
string value that follows into \PB{\\{out}} (at most \PB{$\\{max\_len}-\T{1}$}
bytes), and
advances \PB{${*}\\{pos}$} past it.  Returns 1 on success, 0 if the key is
absent.

\Y\B\4\X16:JSON scanner\X${}\E{}$\6
\1\1\&{static} \&{int} \\{scan\_string}(\&{const} \&{char} ${}{*}\\{key},\39{}$%
\&{const} \&{char} ${}{*}{*}\\{pos},\39{}$\&{char} ${}{*}\\{out},\39{}$\&{size%
\_t} \\{max\_len})\2\2\6
${}\{{}$\1\6
\&{char} \\{needle}[\T{128}];\7
${}\\{snprintf}(\\{needle},\39{}$\&{sizeof} \\{needle}${},\39\.{"\\"\%s\\""},%
\39\\{key});{}$\7
\&{const} \&{char} ${}{*}\|p\K\\{strstr}({*}\\{pos},\39\\{needle});{}$\7
\&{if} ${}(\R\|p){}$\1\5
\&{return} \T{0};\2\6
${}\|p\MRL{+{\K}}\\{strlen}(\\{needle});{}$\6
\&{while} ${}({*}\|p\E\.{'\ '}\V{*}\|p\E\.{':'}){}$\1\5
${}\|p\PP;{}$\2\6
\&{if} ${}({*}\|p\I\.{'"'}){}$\1\5
\&{return} \T{0};\2\6
${}\|p\PP;{}$\7
\&{const} \&{char} ${}{*}\|e\K\\{strchr}(\|p,\39\.{'"'});{}$\7
\&{if} ${}(\R\|e){}$\1\5
\&{return} \T{0};\2\7
\&{size\_t} \\{len}${}\K(\&{size\_t})(\|e-\|p);{}$\7
\&{if} ${}(\\{len}\G\\{max\_len}){}$\1\5
${}\\{len}\K\\{max\_len}-\T{1};{}$\2\6
${}\\{memcpy}(\\{out},\39\|p,\39\\{len});{}$\6
${}\\{out}[\\{len}]\K\.{'\\0'};{}$\6
${}{*}\\{pos}\K\|e+\T{1};{}$\6
\&{return} \T{1};\6
\4${}\}{}$\2\par
\A17.
\U4.\fi

\M{17}\PB{\\{scan\_long}} is the numeric counterpart: it finds \PB{\.{"key"}},
converts
the digits that follow to a \PB{\&{long}}, stores the result in \PB{${*}%
\\{out}$}, and
advances \PB{${*}\\{pos}$}.

\Y\B\4\X16:JSON scanner\X${}\mathrel+\E{}$\6
\1\1\&{static} \&{int} \\{scan\_long}(\&{const} \&{char} ${}{*}\\{key},\39{}$%
\&{const} \&{char} ${}{*}{*}\\{pos},\39{}$\&{long} ${}{*}\\{out})\2\2{}$\6
${}\{{}$\1\6
\&{char} \\{needle}[\T{128}];\7
${}\\{snprintf}(\\{needle},\39{}$\&{sizeof} \\{needle}${},\39\.{"\\"\%s\\""},%
\39\\{key});{}$\7
\&{const} \&{char} ${}{*}\|p\K\\{strstr}({*}\\{pos},\39\\{needle});{}$\7
\&{if} ${}(\R\|p){}$\1\5
\&{return} \T{0};\2\6
${}\|p\MRL{+{\K}}\\{strlen}(\\{needle});{}$\6
\&{while} ${}({*}\|p\E\.{'\ '}\V{*}\|p\E\.{':'}){}$\1\5
${}\|p\PP;{}$\2\7
\&{char} ${}{*}\\{end};{}$\7
${}{*}\\{out}\K\\{strtol}(\|p,\39{\AND}\\{end},\39\T{10});{}$\6
\&{if} ${}(\\{end}\E\|p){}$\1\5
\&{return} \T{0};\2\6
${}{*}\\{pos}\K\\{end};{}$\6
\&{return} \T{1};\6
\4${}\}{}$\2\par
\fi

\N{1}{18}Person lookup.

\fi

\M{19}\PB{\\{lookup\_person\_key}} posts a search for \PB{\\{search\_query}}
and scans
the result rows for one whose full name contains \PB{\\{match\_substr}} (after
lower-casing).  It returns a heap-allocated decimal \PB{\\{PersonKey}} string
and, optionally, the swimmer's full name in \PB{${*}\\{out\_name}$}.  Returns %
\PB{$\NULL$}
on failure.

\Y\B\4\X19:Person lookup\X${}\E{}$\6
\1\1\&{static} \&{char} ${}{*}{}$\\{lookup\_person\_key}(\&{const} \&{char}
${}{*}\\{search\_query},\39{}$\&{const} \&{char} ${}{*}\\{match\_substr},\39{}$%
\&{const} \&{char} ${}{*}{*}\\{out\_name})\2\2{}$\6
${}\{{}$\1\6
\&{char} \\{url}[\T{512}];\7
${}\\{snprintf}(\\{url},\39{}$\&{sizeof} \\{url}${},\39\.{"%
\%s/aPublicIAAaPerso}\)\.{nIAAaSearch/jaql"},\39\.{SISENSE\_API});{}$\7
\&{char} \\{body}[\T{1024}];\7
${}\\{snprintf}(\\{body},\39{}$\&{sizeof} \\{body}${},\39\.{"\{"}\.{"%
\\"datasource\\":\{"}\.{"\\"title\\":\\"Public\ }\)\.{Person\ Search\\","}\.{"%
\\"fullname\\":\\"Loca}\)\.{lHost/Public\ Person\ }\)\.{Search\\"\},"}\.{"%
\\"metadata\\":["}\.{"\{\\"jaql\\":\{\\"table\\}\)\.{":\\"Persons\\",\\"colu}\)%
\.{mn\\":\\"FullName\\","}\.{"\\"dim\\":\\"[Persons.}\)\.{FullName]\\",%
\\"dataty}\)\.{pe\\":\\"text\\","}\.{"\\"title\\":\\"Name\\",}\)\.{"}\.{"%
\\"filter\\":\{\\"conta}\)\.{ins\\":\\"\%s\\"\}\}\},"}\.{"\{\\"jaql\\":\{%
\\"table\\}\)\.{":\\"Persons\\",\\"colu}\)\.{mn\\":\\"PersonKey\\","}\.{"\\"dim%
\\":\\"[Persons.}\)\.{PersonKey]\\",\\"datat}\)\.{ype\\":\\"numeric\\","}\.{"%
\\"title\\":\\"PersonK}\)\.{ey\\"\}\},"}\.{"\{\\"jaql\\":\{\\"table\\}\)\.{":%
\\"Persons\\",\\"colu}\)\.{mn\\":\\"ClubName\\","}\.{"\\"dim\\":\\"[Persons.}\)%
\.{ClubName]\\",\\"dataty}\)\.{pe\\":\\"text\\","}\.{"\\"title\\":\\"Club\\"\}}%
\)\.{\}"}\.{"],\\"count\\":100,\\"o}\)\.{ffset\\":0\}"},\39\\{search%
\_query});{}$\7
\&{char} ${}{*}\\{resp}\K\\{post\_json}(\\{url},\39\\{body});{}$\7
\&{if} ${}(\R\\{resp}){}$\5
${}\{{}$\1\6
${}\\{fputs}(\.{"Error:\ person\ looku}\)\.{p\ request\ failed\\n"},\39%
\\{stderr});{}$\6
\&{return} ${}\NULL;{}$\6
\4${}\}{}$\2\7
\&{char} ${}{*}\\{key}\K\NULL;{}$\6
\&{char} ${}{*}\\{name}\K\NULL;{}$\6
\&{const} \&{char} ${}{*}\|p\K\\{strstr}(\\{resp},\39\.{"\\"values\\""});{}$\7
\&{while} (\|p)\5
${}\{{}$\1\6
\&{char} \\{full\_name}[\T{256}];\7
\&{if} ${}(\R\\{scan\_string}(\.{"text"},\39{\AND}\|p,\39\\{full\_name},\39{}$%
\&{sizeof} \\{full\_name}))\1\5
\&{break};\C{ Lower-case the name and check for the match substring. }\2\7
\&{char} \\{lower}[\T{256}];\6
\&{size\_t} \\{flen}${}\K\\{strlen}(\\{full\_name});{}$\7
\&{for} (\&{size\_t} \|i${}\K\T{0};{}$ ${}\|i\Z\\{flen};{}$ ${}\|i\PP){}$\1\5
${}\\{lower}[\|i]\K(\&{char})(\\{full\_name}[\|i]\G\.{'A'}\W\\{full\_name}[\|i]%
\Z\.{'Z'}\?\\{full\_name}[\|i]+\T{32}:\\{full\_name}[\|i]);{}$\2\6
\&{if} ${}(\\{strstr}(\\{lower},\39\\{match\_substr})){}$\5
${}\{{}$\1\6
\&{long} \\{pk};\7
\&{if} ${}(\\{scan\_long}(\.{"data"},\39{\AND}\|p,\39{\AND}\\{pk})){}$\5
${}\{{}$\1\6
\&{char} \\{tmp}[\T{32}];\7
${}\\{snprintf}(\\{tmp},\39{}$\&{sizeof} \\{tmp}${},\39\.{"\%ld"},\39%
\\{pk});{}$\6
${}\\{key}\K\\{strdup}(\\{tmp});{}$\6
${}\\{name}\K\\{strdup}(\\{full\_name});{}$\6
\4${}\}{}$\2\6
\&{break};\6
\4${}\}{}$\2\6
\4${}\}{}$\2\6
\\{free}(\\{resp});\6
\&{if} ${}(\R\\{key}){}$\1\5
${}\\{fprintf}(\\{stderr},\39\.{"Error:\ swimmer\ \\"\%s}\)\.{\\"\ not\ found%
\\n"},\39\\{search\_query});{}$\2\6
\&{if} (\\{out\_name})\1\5
${}{*}\\{out\_name}\K\\{name};{}$\2\6
\&{else}\1\5
\\{free}(\\{name});\2\6
\&{return} \\{key};\6
\4${}\}{}$\2\par
\U4.\fi

\N{1}{20}Times fetch.

\fi

\M{21}\PB{\\{insertion\_sort}} sorts \PB{\|n} \PB{\&{TimeRow}} records in-place
by \PB{\\{sort\_key}}
ascending (smallest = fastest time).  A swimmer's career rarely exceeds
a few dozen entries per event, so $O(n^2)$ is entirely adequate.

\Y\B\4\X21:Times fetch\X${}\E{}$\6
\1\1\&{static} \&{void} \\{insertion\_sort}(\&{TimeRow} ${}{*}\\{rows},\39{}$%
\&{int} \|n)\2\2\6
${}\{{}$\1\6
\&{for} (\&{int} \|i${}\K\T{1};{}$ ${}\|i<\|n;{}$ ${}\|i\PP){}$\5
${}\{{}$\1\6
\&{TimeRow} \\{tmp}${}\K\\{rows}[\|i];{}$\6
\&{int} \|j${}\K\|i-\T{1};{}$\7
\&{while} ${}(\|j\G\T{0}\W\\{rows}[\|j].\\{sort\_key}>\\{tmp}.\\{sort\_key}){}$%
\5
${}\{{}$\1\6
${}\\{rows}[\|j+\T{1}]\K\\{rows}[\|j];{}$\6
${}\|j\MM;{}$\6
\4${}\}{}$\2\6
${}\\{rows}[\|j+\T{1}]\K\\{tmp};{}$\6
\4${}\}{}$\2\6
\4${}\}{}$\2\par
\As22\ET23.
\U4.\fi

\M{22}\PB{\\{format\_date}} converts an 8-digit \.{YYYYMMDD} string (such as
\.{20230305}) to \.{YYYY-MM-DD} format in \PB{\\{out}} (which must be at
least 11 bytes).

\Y\B\4\X21:Times fetch\X${}\mathrel+\E{}$\6
\1\1\&{static} \&{void} \\{format\_date}(\&{const} \&{char} ${}{*}\\{ymd8},%
\39{}$\&{char} ${}{*}\\{out})\2\2{}$\6
${}\{{}$\1\6
\&{if} ${}(\\{strlen}(\\{ymd8})\E\T{8}){}$\5
${}\{{}$\1\6
${}\\{out}[\T{0}]\K\\{ymd8}[\T{0}];{}$\6
${}\\{out}[\T{1}]\K\\{ymd8}[\T{1}];{}$\6
${}\\{out}[\T{2}]\K\\{ymd8}[\T{2}];{}$\6
${}\\{out}[\T{3}]\K\\{ymd8}[\T{3}];{}$\6
${}\\{out}[\T{4}]\K\.{'-'};{}$\6
${}\\{out}[\T{5}]\K\\{ymd8}[\T{4}];{}$\6
${}\\{out}[\T{6}]\K\\{ymd8}[\T{5}];{}$\6
${}\\{out}[\T{7}]\K\.{'-'};{}$\6
${}\\{out}[\T{8}]\K\\{ymd8}[\T{6}];{}$\6
${}\\{out}[\T{9}]\K\\{ymd8}[\T{7}];{}$\6
${}\\{out}[\T{10}]\K\.{'\\0'};{}$\6
\4${}\}{}$\2\6
\&{else}\5
${}\{{}$\1\6
${}\\{strncpy}(\\{out},\39\\{ymd8},\39\T{10});{}$\6
${}\\{out}[\T{10}]\K\.{'\\0'};{}$\6
\4${}\}{}$\2\6
\4${}\}{}$\2\par
\fi

\M{23}\PB{\\{fetch\_times}} retrieves all SCY times for \PB{\\{person\_key}} in
\PB{\\{event\_code}}, sorts them fastest-first, and prints them.  Each JAQL
result row delivers five cells in order: time, sort key, meet name,
swim date (as \.{YYYYMMDD}), and motivational standard type.
The sort key (a decimal float string such as \.{"1010007029.00"}) is
converted to \PB{\&{double}} with \PB{\\{strtod}}.  The date is reformatted
from
\.{YYYYMMDD} to \.{YYYY-MM-DD} by \PB{\\{format\_date}}.

\Y\B\4\X21:Times fetch\X${}\mathrel+\E{}$\6
\1\1\&{static} \&{void} \\{fetch\_times}(\&{const} \&{char} ${}{*}\\{person%
\_key},\39{}$\&{const} \&{char} ${}{*}\\{event\_code})\2\2{}$\6
${}\{{}$\1\6
\&{char} \\{url}[\T{512}];\7
${}\\{snprintf}(\\{url},\39{}$\&{sizeof} \\{url}${},\39\.{"%
\%s/aUSAIAAaSwimming}\)\.{IAAaTimesIAAaElastic}\)\.{ube/jaql"},\39\.{SISENSE%
\_API});{}$\7
\&{char} \\{body}[\T{2048}];\7
${}\\{snprintf}(\\{body},\39{}$\&{sizeof} \\{body}${},\39\.{"\{"}\.{"%
\\"datasource\\":\{"}\.{"\\"title\\":\\"USA\ Swi}\)\.{mming\ Times\ Elasticu}\)%
\.{be\\","}\.{"\\"fullname\\":\\"Loca}\)\.{lHost/USA\ Swimming\ T}\)\.{imes\
Elasticube\\"\},"}\.{"\\"metadata\\":["}\.{"\{\\"jaql\\":\{\\"table\\}\)\.{":%
\\"UsasSwimTime\\",\\}\)\.{"column\\":\\"PersonKe}\)\.{y\\","}\.{"\\"dim\\":%
\\"[UsasSwim}\)\.{Time.PersonKey]\\",\\"}\)\.{datatype\\":\\"numeric}\)\.{%
\\","}\.{"\\"title\\":\\"PersonK}\)\.{ey\\","}\.{"\\"filter\\":\{\\"equal}\)%
\.{s\\":\%s\}\},\\"panel\\":\\}\)\.{"scope\\"\},"}\.{"\{\\"jaql\\":\{\\"table%
\\}\)\.{":\\"SwimEvent\\",\\"co}\)\.{lumn\\":\\"EventCode\\"}\)\.{,"}\.{"\\"dim%
\\":\\"[SwimEven}\)\.{t.EventCode]\\",\\"dat}\)\.{atype\\":\\"text\\","}\.{"%
\\"title\\":\\"Event\\"}\)\.{,"}\.{"\\"filter\\":\{\\"equal}\)\.{s\\":\\"\%s\\"%
\}\},\\"panel}\)\.{\\":\\"scope\\"\},"}\.{"\{\\"jaql\\":\{\\"table\\}\)\.{":%
\\"UsasSwimTime\\","}\.{"\\"column\\":\\"SwimTi}\)\.{meFormatted\\","}\.{"%
\\"dim\\":\\"[UsasSwim}\)\.{Time.SwimTimeFormatt}\)\.{ed]\\","}\.{"\\"datatype%
\\":\\"text}\)\.{\\",\\"title\\":\\"Time\\}\)\.{"\}\},"}\.{"\{\\"jaql\\":\{%
\\"table\\}\)\.{":\\"UsasSwimTime\\",\\}\)\.{"column\\":\\"SortKey\\}\)\.{","}%
\.{"\\"dim\\":\\"[UsasSwim}\)\.{Time.SortKey]\\",\\"da}\)\.{tatype\\":%
\\"numeric\\"}\)\.{,"}\.{"\\"title\\":\\"SortKey}\)\.{\\"\}\},"}\.{"\{\\"jaql%
\\":\{\\"table\\}\)\.{":\\"Meet\\",\\"column\\}\)\.{":\\"MeetName\\","}\.{"%
\\"dim\\":\\"[Meet.Mee}\)\.{tName]\\",\\"datatype\\}\)\.{":\\"text\\","}\.{"%
\\"title\\":\\"Meet\\"\}}\)\.{\},"}\.{"\{\\"jaql\\":\{\\"table\\}\)\.{":%
\\"UsasSwimTime\\","}\.{"\\"column\\":\\"Season}\)\.{CalendarKey\\","}\.{"%
\\"dim\\":\\"[UsasSwim}\)\.{Time.SeasonCalendarK}\)\.{ey]\\","}\.{"\\"datatype%
\\":\\"nume}\)\.{ric\\",\\"title\\":\\"Da}\)\.{te\\"\}\},"}\.{"\{\\"jaql\\":\{%
\\"table\\}\)\.{":\\"TimeStandard\\","}\.{"\\"column\\":\\"Standa}\)\.{rdType%
\\","}\.{"\\"dim\\":\\"[TimeStan}\)\.{dard.StandardType]\\"}\)\.{,"}\.{"%
\\"datatype\\":\\"text}\)\.{\\",\\"title\\":\\"Stand}\)\.{ard\\"\}\}"}\.{"],%
\\"count\\":100,\\"o}\)\.{ffset\\":0\}"},\39\\{person\_key},\39\\{event%
\_code});{}$\7
\&{char} ${}{*}\\{resp}\K\\{post\_json}(\\{url},\39\\{body});{}$\7
\&{if} ${}(\R\\{resp}){}$\5
${}\{{}$\1\6
${}\\{fprintf}(\\{stderr},\39\.{"Error:\ times\ reques}\)\.{t\ failed\ for\ \%s%
\\n"},\39\\{event\_code});{}$\6
\&{return};\6
\4${}\}{}$\2\7
\&{TimeRow} \\{rows}[\.{MAX\_TIMES}];\6
\&{int} \\{nrows}${}\K\T{0};{}$\6
\&{const} \&{char} ${}{*}\|p\K\\{strstr}(\\{resp},\39\.{"\\"values\\""}){}$;\C{
Each row has five cells scanned in order via "text":        0 = swim time
string, 1 = sort key as decimal float string,        2 = meet name, 3 = date as
YYYYMMDD, 4 = motivational standard. }\7
\&{while} ${}(\|p\W\\{nrows}<\.{MAX\_TIMES}){}$\5
${}\{{}$\1\6
\&{char} \\{time\_str}[\T{32}];\7
\&{if} ${}(\R\\{scan\_string}(\.{"text"},\39{\AND}\|p,\39\\{time\_str},\39{}$%
\&{sizeof} \\{time\_str}))\1\5
\&{break};\2\6
\&{if} ${}(\\{time\_str}[\T{0}]\E\.{'\\0'}){}$\1\5
\&{break};\2\7
\&{char} \\{sort\_str}[\T{64}];\7
\&{if} ${}(\R\\{scan\_string}(\.{"text"},\39{\AND}\|p,\39\\{sort\_str},\39{}$%
\&{sizeof} \\{sort\_str}))\1\5
\&{break};\2\7
\&{char} \\{meet\_str}[\T{256}];\7
\&{if} ${}(\R\\{scan\_string}(\.{"text"},\39{\AND}\|p,\39\\{meet\_str},\39{}$%
\&{sizeof} \\{meet\_str}))\1\5
\&{break};\2\7
\&{char} \\{date\_str}[\T{16}];\7
\&{if} ${}(\R\\{scan\_string}(\.{"text"},\39{\AND}\|p,\39\\{date\_str},\39{}$%
\&{sizeof} \\{date\_str}))\1\5
\&{break};\2\7
\&{char} \\{std\_str}[\T{48}];\7
\&{if} ${}(\R\\{scan\_string}(\.{"text"},\39{\AND}\|p,\39\\{std\_str},\39{}$%
\&{sizeof} \\{std\_str}))\1\5
\&{break};\2\6
${}\\{rows}[\\{nrows}].\\{sort\_key}\K\\{strtod}(\\{sort\_str},\39\NULL);{}$\6
${}\\{strncpy}(\\{rows}[\\{nrows}].\\{time},\39\\{time\_str},\39{}$\&{sizeof} %
\\{rows}[\\{nrows}]${}.\\{time}-\T{1});{}$\6
${}\\{rows}[\\{nrows}].\\{time}{}$[\&{sizeof} \\{rows}[\\{nrows}]${}.\\{time}-%
\T{1}]\K\.{'\\0'};{}$\6
${}\\{format\_date}(\\{date\_str},\39\\{rows}[\\{nrows}].\\{date});{}$\6
${}\\{strncpy}(\\{rows}[\\{nrows}].\\{standard},\39\\{std\_str},\39{}$%
\&{sizeof} \\{rows}[\\{nrows}]${}.\\{standard}-\T{1});{}$\6
${}\\{rows}[\\{nrows}].\\{standard}{}$[\&{sizeof} \\{rows}[\\{nrows}]${}.%
\\{standard}-\T{1}]\K\.{'\\0'};{}$\6
${}\\{strncpy}(\\{rows}[\\{nrows}].\\{meet},\39\\{meet\_str},\39{}$\&{sizeof} %
\\{rows}[\\{nrows}]${}.\\{meet}-\T{1});{}$\6
${}\\{rows}[\\{nrows}].\\{meet}{}$[\&{sizeof} \\{rows}[\\{nrows}]${}.\\{meet}-%
\T{1}]\K\.{'\\0'};{}$\6
${}\\{nrows}\PP;{}$\6
\4${}\}{}$\2\6
\\{free}(\\{resp});\6
${}\\{insertion\_sort}(\\{rows},\39\\{nrows});{}$\6
${}\\{printf}(\.{"\%s\ ---\ all\ times\ (f}\)\.{astest\ first):\\n"},\39%
\\{event\_code});{}$\6
${}\\{printf}(\.{"\%-12s\ \ \%-10s\ \ \%-13s}\)\.{\ \ \%s\\n"},\39\.{"Time"},%
\39\.{"Date"},\39\.{"Standard"},\39\.{"Meet"});{}$\6
${}\\{printf}(\.{"\%-12s\ \ \%-10s\ \ \%-13s}\)\.{\ \ \%s\\n"},\39%
\.{"------------"},\39\.{"----------"},\39\.{"-------------"},\39%
\.{"----"});{}$\6
\&{for} (\&{int} \|i${}\K\T{0};{}$ ${}\|i<\\{nrows};{}$ ${}\|i\PP){}$\1\5
${}\\{printf}(\.{"\%-12s\ \ \%-10s\ \ \%-13s}\)\.{\ \ \%s\\n"},\39\\{rows}[%
\|i].\\{time},\39\\{rows}[\|i].\\{date},\39\\{rows}[\|i].\\{standard},\39%
\\{rows}[\|i].\\{meet});{}$\2\6
\&{if} ${}(\\{nrows}\E\T{0}){}$\1\5
\\{printf}(\.{"(no\ times\ found)\\n"});\2\6
\\{putchar}(\.{'\\n'});\6
\4${}\}{}$\2\par
\fi

\N{1}{24}Main program.

\fi

\M{25}We define a small \PB{\\{Swimmer}} record to hold the person-search
parameters for each swimmer.  Each entry supplies a \PB{\\{search\_query}}
string
sent to the database (ideally distinctive enough to return a small set)
and a \PB{\\{match\_substr}} (lower-cased) used to identify the correct row.
Kalea's entry searches ``Benavente'' because her registered name is
``Kalea Rose Benavente''; the simple two-word query ``Kalea Benavente''
returns no results since the API requires an exact substring match.

\Y\B\4\X25:Main function\X${}\E{}$\6
\&{typedef} \&{struct} ${}\{{}$\1\6
\&{const} \&{char} ${}{*}\\{search\_query}{}$;\C{ Name string to search for }\6
\&{const} \&{char} ${}{*}\\{match\_substr}{}$;\C{ Lower-case substring to match
}\2\6
${}\}{}$ \&{Swimmer};\6
\&{static} \&{const} \&{Swimmer} \.{SWIMMERS}[\,]${}\K\{\{\.{"Julianna\
Evans"},\39\.{"stella"}\},{}$\C{ finds Stella Julianna Evans }\6
${}\{\.{"Benavente"},\39\.{"kalea"}\}{}$\C{ finds Kalea Rose Benavente }\6
${}\};{}$\6
\8\#\&{define} \.{NUM\_SWIMMERS}\5((\&{int})(\&{sizeof} \.{SWIMMERS}${}/{}$%
\&{sizeof} \.{SWIMMERS}[\T{0}]))\par
\A26.
\U4.\fi

\M{26}We initialise the global \.{libcurl} state, then for each swimmer
resolve her \PB{\\{PersonKey}}, print a header, and iterate over all twelve
events, then release all resources before exit.

\Y\B\4\X25:Main function\X${}\mathrel+\E{}$\6
\1\1\&{int} \\{main}(\&{void})\2\2\6
${}\{{}$\1\6
\\{curl\_global\_init}(\.{CURL\_GLOBAL\_DEFAULT});\6
\&{for} (\&{int} \|s${}\K\T{0};{}$ ${}\|s<\.{NUM\_SWIMMERS};{}$ ${}\|s\PP){}$\5
${}\{{}$\1\6
\&{const} \&{char} ${}{*}\\{name}\K\NULL;{}$\6
\&{char} ${}{*}\\{key}\K\\{lookup\_person\_key}(\.{SWIMMERS}[\|s].\\{search%
\_query},\39\.{SWIMMERS}[\|s].\\{match\_substr},\39{\AND}\\{name});{}$\7
\&{if} ${}(\R\\{key}){}$\5
${}\{{}$\1\6
\\{curl\_global\_cleanup}(\,);\6
\&{return} \T{1};\6
\4${}\}{}$\2\6
${}\\{printf}(\.{"Swimmer:\ \%s\ \ (Perso}\)\.{nKey:\ \%s)\\n\\n"},\39\\{name}%
\?\\{name}:\.{"(unknown)"},\39\\{key});{}$\6
\&{for} (\&{int} \|i${}\K\T{0};{}$ ${}\|i<\.{NUM\_EVENTS};{}$ ${}\|i\PP){}$\1\5
${}\\{fetch\_times}(\\{key},\39\.{EVENTS}[\|i]);{}$\2\6
\\{free}(\\{key});\6
\\{free}((\&{void} ${}{*}){}$ \\{name});\6
\\{printf}(\.{"\\n"});\6
\4${}\}{}$\2\6
\\{curl\_global\_cleanup}(\,);\6
\&{return} \T{0};\6
\4${}\}{}$\2\par
\fi

\N{1}{27}Glossary.

The following terms and interfaces appear throughout this program.

\fi

\M{28}
\def\gitem#1{\medskip\noindent{\bf #1.}\enspace\ignorespaces}
\def\sig#1{\par\noindent\quad{\tt #1}\par\noindent}

\gitem{JAQL (JSON Analytics Query Language)}
The query language used by the Sisense Analytic Engine to describe data
requests.  A JAQL query is a JSON object posted to the endpoint
\.{/api/datasources/\{name\}/jaql}.  Its \.{metadata} array specifies
the columns (dimensions or measures) to retrieve; \.{filter} objects
restrict row membership; and \.{count}/\.{offset} control pagination.
The Sisense back-end translates JAQL to an internal columnar query,
executes it against the ElastiCube, and returns results in a
\.{values} array whose rows are parallel to the \.{metadata} array.

\gitem{Sisense Analytic Engine (ElastiCube)}
A columnar, in-memory analytics database developed by Sisense.
USA Swimming hosts a Sisense instance at
\.{usaswimming.sisense.com} that powers the public data hub at
\.{data.usaswimming.org}.  The engine stores swim-time, meet, person,
and time-standard data in compressed column stores called
{\it ElastiCubes}; this program queries two of them---``Public Person
Search'' and ``USA Swimming Times Elasticube''---via their JAQL
endpoints, authenticating with a bearer token in the
\.{Authorization} HTTP header.

\fi

\M{29}{\bf Sisense JAQL API Calls.}
All requests are HTTP POST to %
\.{https://usaswimming.sisense.com/api/datasources/\{ds\}/jaql}
with headers \.{Content-Type: application/json} and
\.{Authorization: Bearer \{token\}}.
The JSON request body has the following top-level fields:

\medskip
\item{$\bullet$} \.{datasource} (object).
Identifies the ElastiCube to query.
\par\noindent Fields:
\itemitem{--} \.{title} (string): human-readable name,
e.g.\ \.{"Public Person Search"}.
\itemitem{--} \.{fullname} (string): internal path used by the server,
e.g.\ \.{"LocalHost/Public Person Search"}.

\item{$\bullet$} \.{metadata} (array of objects).
Each element describes one column of the result.
\par\noindent Per-element fields:
\itemitem{--} \.{jaql.table} (string): the ElastiCube table name
(e.g.\ \.{"Persons"}, \.{"UsasSwimTime"}, \.{"Meet"},
\.{"SwimEvent"}, \.{"TimeStandard"}).
\itemitem{--} \.{jaql.column} (string): the column name within the table
(e.g.\ \.{"FullName"}, \.{"PersonKey"}, \.{"SwimTimeFormatted"},
\.{"SortKey"}, \.{"MeetName"}, \.{"SeasonCalendarKey"},
\.{"EventCode"}, \.{"StandardType"}).
\itemitem{--} \.{jaql.dim} (string): the bracketed dimension path
\.{"[Table.Column]"} used by the query engine.
\itemitem{--} \.{jaql.datatype} (string): \.{"text"} or \.{"numeric"}.
\itemitem{--} \.{jaql.title} (string): label for the column in the response.
\itemitem{--} \.{jaql.filter} (object, optional): restricts rows.
This program uses two filter types:
\.{\{"contains": "..."\}} for substring matches on text columns
(person search) and \.{\{"equals": value\}} for exact matches
({\tt PersonKey} and {\tt EventCode} scope filters).
\itemitem{--} \.{panel} (string, optional): when set to \.{"scope"}
the column acts as a {\it scope filter}---it narrows the result set
without appearing as an output column.

\item{$\bullet$} \.{count} (integer).
Maximum number of rows to return.
This program uses 10 for quick existence checks and 100 for full
result sets.

\item{$\bullet$} \.{offset} (integer).
Zero-based row offset for pagination.

\medskip\noindent
The response body is a JSON object.  Relevant fields:

\medskip
\item{$\bullet$} \.{values} (array of arrays).
Each inner array is one result row; its elements correspond
positionally to the non-scope entries in \.{metadata}.
Each element is an object with two keys: \.{data} (the raw value,
number or string) and \.{text} (the formatted string representation).
This program always reads the \.{text} field.

\item{$\bullet$} \.{error} (boolean, present on failure).
When true, \.{details} contains an error message.

\medskip\noindent
{\bf Person search call.}
Endpoint: \.{.../aPublicIAAaPersonIAAaSearch/jaql}.
Queries the \.{Persons} table with a \.{contains} filter on
\.{FullName}.  Returns \.{FullName} (text) and \.{PersonKey} (numeric)
columns.  This program sends \.{count:100} and scans rows until it
finds one whose lower-cased name contains the match substring.

\medskip\noindent
{\bf Times query call.}
Endpoint: \.{.../aUSAIAAaSwimmingIAAaTimesIAAaElasticube/jaql}.
Uses two scope filters (\.{PersonKey} equals and \.{EventCode} equals)
and retrieves five output columns in order:
\.{SwimTimeFormatted} (formatted time string),
\.{SortKey} (numeric sort key as decimal string, smaller = faster),
\.{MeetName} (meet name),
\.{SeasonCalendarKey} (swim date as \.{YYYYMMDD} integer), and
\.{StandardType} from the \.{TimeStandard} table
(motivational level: \.{"B"}, \.{"BB"}, \.{"A"}, \.{"Slower Than B"}, etc.).

\fi

\M{30}{\bf libcurl API Calls.}
This program uses the libcurl ``easy'' interface for synchronous HTTP.
All functions return a \PB{\\{CURLcode}} (zero = \.{CURLE\_OK}) except where
noted.

\medskip
\item{$\bullet$} {\tt curl\_global\_init(flags)}.
\par\noindent Parameter: {\tt flags} ({\tt long}) ---
a bitmask of subsystems to initialise.
This program passes \.{CURL\_GLOBAL\_DEFAULT}, which enables SSL
and the Windows socket layer on that platform.
Must be called once before any other libcurl function.
Returns a {\tt CURLcode}; this program ignores the return value
because failure is treated as fatal by the subsequent easy calls.

\item{$\bullet$} {\tt curl\_global\_cleanup(void)}.
Releases all resources allocated by \PB{\\{curl\_global\_init}}.
Must be called once after all easy handles have been cleaned up.
Returns nothing.

\item{$\bullet$} {\tt curl\_easy\_init(void)}.
Allocates and returns a new easy handle (a \.{CURL *}).
Returns \.{NULL} on failure.
Each call to \PB{\\{post\_json}} creates its own handle and destroys it
before returning, so handles are never shared between requests.

\item{$\bullet$} {\tt curl\_easy\_setopt(handle, option, value)}.
\par\noindent Parameters:
\itemitem{--} {\tt handle} ({\tt CURL *}): the easy handle.
\itemitem{--} {\tt option} ({\tt CURLoption}): a constant selecting
the behaviour to configure.  Options used here:
\itemitem{} \.{CURLOPT\_URL} ({\tt char *}) --- the request URL.
\itemitem{} \.{CURLOPT\_HTTPHEADER} ({\tt struct curl\_slist *}) ---
linked list of extra HTTP headers
(\.{Content-Type} and \.{Authorization}).
\itemitem{} \.{CURLOPT\_POSTFIELDS} ({\tt char *}) ---
the POST body; setting this also switches the method to POST.
\itemitem{} \.{CURLOPT\_WRITEFUNCTION} (function pointer) ---
callback invoked for each response chunk; signature
{\tt size\_t cb(void*,size\_t,size\_t,void*)}.
\itemitem{} \.{CURLOPT\_WRITEDATA} ({\tt void *}) ---
the user-data pointer passed as the fourth argument to the
write callback; here a pointer to the \PB{\&{Buffer}} accumulator.
\itemitem{--} {\tt value}: type depends on {\tt option} (see above).

\item{$\bullet$} {\tt curl\_easy\_perform(handle)}.
\par\noindent Parameter: {\tt handle} ({\tt CURL *}).
Executes the configured request synchronously, invoking the write
callback for each received chunk.
Returns \.{CURLE\_OK} on success or a non-zero error code; on failure
\PB{\\{post\_json}} frees the partial buffer and returns \.{NULL}.

\item{$\bullet$} {\tt curl\_easy\_cleanup(handle)}.
\par\noindent Parameter: {\tt handle} ({\tt CURL *}).
Releases all resources associated with the handle.
The handle must not be used after this call.

\item{$\bullet$} {\tt curl\_slist\_append(list, string)}.
\par\noindent Parameters:
\itemitem{--} {\tt list} ({\tt struct curl\_slist *}):
existing list head, or \.{NULL} to start a new list.
\itemitem{--} {\tt string} ({\tt const char *}): the string to append.
Returns the new list head, or \.{NULL} on allocation failure.
Used to build the two-header list
(\.{Content-Type} then \.{Authorization}).

\item{$\bullet$} {\tt curl\_slist\_free\_all(list)}.
\par\noindent Parameter: {\tt list} ({\tt struct curl\_slist *}).
Frees every node in the linked list.
Called immediately after \PB{\\{curl\_easy\_perform}} so the headers are
released before the handle.

\fi

\M{31}{\bf POSIX System Calls and Library Functions.}
The following identifiers from the POSIX.1-2008 standard are used
directly in this program.  Each entry gives the C~signature, a
description of each parameter, and a note on how the program uses it.

\medskip
\item{$\bullet$} {\tt int fprintf(FILE *stream, const char *fmt, ...)}.
\par\noindent Parameters:
\itemitem{--} {\tt stream}: destination file (\.{stderr} here).
\itemitem{--} {\tt fmt}: printf-style format string.
\itemitem{--} {\tt ...}: values substituted into {\tt fmt}.
Writes formatted output to {\tt stream}; returns the character count
or a negative value on error.
Used to report lookup and HTTP failures to standard error.

\item{$\bullet$} {\tt int fputs(const char *s, FILE *stream)}.
\par\noindent Parameters:
\itemitem{--} {\tt s}: null-terminated string to write.
\itemitem{--} {\tt stream}: destination file (\.{stderr} here).
Writes {\tt s} without a trailing newline; returns non-negative on
success or \.{EOF} on error.
Used for fixed error messages where no formatting is needed.

\item{$\bullet$} {\tt void free(void *ptr)}.
\par\noindent Parameter:
\itemitem{--} {\tt ptr}: pointer to a heap block, or \.{NULL}
(in which case nothing happens).
Releases the block back to the heap.
Called on every heap string (response buffers, duplicated names and
keys) when they are no longer needed.

\item{$\bullet$} {\tt void *memcpy(void *dst, const void *src, size\_t n)}.
\par\noindent Parameters:
\itemitem{--} {\tt dst}: destination address.
\itemitem{--} {\tt src}: source address.
\itemitem{--} {\tt n}: number of bytes to copy.
Copies exactly {\tt n} bytes from {\tt src} to {\tt dst};
regions must not overlap.
Returns {\tt dst}.
Used in \PB{\\{write\_cb}} to append each network chunk to the buffer,
and in \PB{\\{scan\_string}} to copy a JSON string value.

\item{$\bullet$} {\tt int printf(const char *fmt, ...)}.
\par\noindent Parameters:
\itemitem{--} {\tt fmt}: printf-style format string.
\itemitem{--} {\tt ...}: values substituted into {\tt fmt}.
Writes formatted output to standard output; returns the character
count or a negative value on error.
Used for all swimmer and time-table output.

\item{$\bullet$} {\tt int putchar(int c)}.
\par\noindent Parameter:
\itemitem{--} {\tt c}: character value (as {\tt unsigned char} cast
to {\tt int}).
Writes one character to standard output; returns the character
written, or \.{EOF} on error.
Used to emit a blank line (\.{'\char`\\n'}) after each event section.

\item{$\bullet$} {\tt void *realloc(void *ptr, size\_t size)}.
\par\noindent Parameters:
\itemitem{--} {\tt ptr}: existing heap block, or \.{NULL}.
\itemitem{--} {\tt size}: new size in bytes.
Returns a pointer to the resized block (possibly moved), or \.{NULL}
if allocation fails (the original block is unchanged on failure).
When {\tt ptr} is \.{NULL} the call is equivalent to {\tt malloc}.
Used in \PB{\\{write\_cb}} to grow the response buffer incrementally as
libcurl delivers each chunk.

\item{$\bullet$} {\tt int snprintf(char *buf, size\_t n, const char *fmt,
...)}.
\par\noindent Parameters:
\itemitem{--} {\tt buf}: destination character array.
\itemitem{--} {\tt n}: maximum bytes to write, including the null terminator.
\itemitem{--} {\tt fmt}: printf-style format string.
\itemitem{--} {\tt ...}: values substituted into {\tt fmt}.
Writes at most {\tt n}$-1$ formatted characters to {\tt buf} and
always null-terminates.  Returns the number of characters that would
have been written had the buffer been unlimited (so a return value
$\ge${\tt n} signals truncation).
Used to assemble URL strings, JSON bodies, and the bearer-token header.

\item{$\bullet$} {\tt char *strchr(const char *s, int c)}.
\par\noindent Parameters:
\itemitem{--} {\tt s}: string to search.
\itemitem{--} {\tt c}: character to find (compared as {\tt unsigned char}).
Returns a pointer to the first occurrence of {\tt c} in {\tt s},
including the terminator if {\tt c} is \.{'\char`\\0'}, or \.{NULL}.
Used in \PB{\\{scan\_string}} to find the closing double-quote of a JSON
string value.

\item{$\bullet$} {\tt char *strdup(const char *s)}.
\par\noindent Parameter:
\itemitem{--} {\tt s}: null-terminated string to duplicate.
Allocates a new heap block of {\tt strlen(s)+1} bytes, copies
{\tt s} into it, and returns the pointer; returns \.{NULL} on failure.
Used in \PB{\\{lookup\_person\_key}} to persist the swimmer's full name and
PersonKey string across the lifetime of a query.

\item{$\bullet$} {\tt size\_t strlen(const char *s)}.
\par\noindent Parameter:
\itemitem{--} {\tt s}: null-terminated string.
Returns the number of bytes before the null terminator.
Used to compute loop bounds when lower-casing names, to advance
past a search needle in \PB{\\{scan\_string}} and \PB{\\{scan\_long}}, and to
validate the 8-character date string in \PB{\\{format\_date}}.

\item{$\bullet$} {\tt char *strncpy(char *dst, const char *src, size\_t n)}.
\par\noindent Parameters:
\itemitem{--} {\tt dst}: destination array (at least {\tt n} bytes).
\itemitem{--} {\tt src}: source string.
\itemitem{--} {\tt n}: maximum bytes to copy.
Copies up to {\tt n} bytes; if {\tt src} is shorter than {\tt n},
the remainder of {\tt dst} is zero-filled.  If {\tt src} is at least
{\tt n} bytes long, {\tt dst} will {\it not\/} be null-terminated.
Returns {\tt dst}.
This program always writes {\tt dst[sizeof dst - 1] = '\char`\\0'}
after the call to guarantee termination.
Used to copy time, standard, and meet strings into {\tt TimeRow}.

\item{$\bullet$} {\tt char *strstr(const char *hay, const char *needle)}.
\par\noindent Parameters:
\itemitem{--} {\tt hay}: string to search within.
\itemitem{--} {\tt needle}: substring to search for.
Returns a pointer to the first occurrence of {\tt needle} in
{\tt hay}, or \.{NULL} if not found.
The workhorse of the JSON scanner: used to locate key names
(\.{"text"}, \.{"data"}, \.{"values"}) and to advance through
the raw Sisense response text.

\item{$\bullet$} {\tt double strtod(const char *s, char **endptr)}.
\par\noindent Parameters:
\itemitem{--} {\tt s}: string containing a floating-point number.
\itemitem{--} {\tt endptr}: if non-\.{NULL}, receives a pointer to the
first character not consumed by the conversion.
Returns the parsed {\tt double}; sets {\tt *endptr} past the
converted text.
Used to convert the Sisense sort-key string
(e.g.\ \.{"1010007029.00"}) to a {\tt double} for the insertion sort.

\item{$\bullet$} {\tt long strtol(const char *s, char **endptr, int base)}.
\par\noindent Parameters:
\itemitem{--} {\tt s}: string containing an integer.
\itemitem{--} {\tt endptr}: if non-\.{NULL}, receives a pointer to the
first character not consumed.
\itemitem{--} {\tt base}: numeric base (2--36), or 0 for auto-detection
from a \.{0x} or \.{0} prefix.  This program passes 10 (decimal).
Returns the parsed {\tt long}; sets {\tt *endptr} past the converted
text.  Used in \PB{\\{scan\_long}} to read the \.{PersonKey} integer from
the JSON person-search response.

\fi

\N{1}{32}Index.

\fi


\inx
\fin
\con
